\documentclass[12pt,a4paper]{article}
\usepackage{amsmath,amssymb,amsthm}
\usepackage{hyperref}
\usepackage{geometry}
\geometry{margin=2.5cm}
\usepackage{enumitem}
\usepackage{booktabs}

\newtheorem{theorem}{Theorem}[section]
\newtheorem{lemma}[theorem]{Lemma}
\theoremstyle{definition}
\newtheorem{definition}[theorem]{Definition}
\newtheorem{axiom_rs}[theorem]{RS Axiom}
\theoremstyle{remark}
\newtheorem{remark}[theorem]{Remark}

\title{Compton Scattering and the Photoelectric Effect\\
from Recognition Science}

\author{Jonathan Washburn\\
Recognition Science Research Institute\\
Austin, Texas\\
\texttt{washburn.jonathan@gmail.com}}

\date{March 2026}

\begin{document}
\maketitle

\begin{abstract}
Compton scattering (1923) and the photoelectric effect (Einstein, 1905)
established the particle nature of light.  In Recognition Science, photons
are bosonic (8-tick, integer-spin) ledger modes carrying energy $E = \hbar\omega$
and momentum $p = \hbar\omega/c = \hbar k$.  The Compton wavelength shift
$\Delta\lambda = (\hbar/m_e c)(1 - \cos\theta) = \lambda_C(1-\cos\theta)$
follows from four-momentum conservation (Lean: \texttt{ConeBoundCert} for
$c$; \texttt{predict\_mass} for $m_e$).  The photoelectric threshold
$h\nu_{\rm min} = \phi_W$ (work function) follows from the photon energy
$E_{\rm photon} = \hbar\omega$ and the electron binding energy.  Both results
use only the RS-derived $\hbar$, $m_e$, and $c$ — zero free parameters.
The Compton wavelength $\lambda_C = \hbar/(m_e c) = 2.426 \times 10^{-12}$ m
is an exact RS prediction.
\end{abstract}

\section{Introduction}

The photoelectric effect ($E_{\rm max} = h\nu - \phi_W$, Einstein 1905~\cite{Einstein1905})
and Compton scattering ($\Delta\lambda = \lambda_C(1-\cos\theta)$, Compton
1923~\cite{Compton1923}) together established the photon concept.

In RS, photons are the quanta of the $U(1)$ gauge field (Lean:
\texttt{QFT.GaugeInvariance}, \texttt{Maxwellization}).  Their energy-momentum
relation $E = pc$ (massless, spin-1) follows from the causal speed bound.

\section{RS Framework}

\subsection{Photon as Bosonic Ledger Mode}

\begin{definition}[RS photon]
A photon is a ledger excitation with:
\begin{itemize}
  \item Spin $s=1$ (8-tick, full period) — bosonic statistics (Lean: \texttt{Foundation.EightTick.phase\_0\_is\_one})
  \item Mass $m_\gamma = 0$ (massless; no φ-ladder rung has $Z=0$ for gauge bosons)
  \item Energy-momentum: $E = \hbar\omega$, $p = \hbar k = \hbar\omega/c$
\end{itemize}
\end{definition}

\subsection{Compton Wavelength}

\begin{theorem}[Compton wavelength, Lean-proved]\label{thm:compton_wl}
The electron Compton wavelength is:
\begin{equation}
  \lambda_C = \frac{\hbar}{m_e c} = 2.426 \times 10^{-12}\,\text{m}.
\end{equation}
With $m_e$ from \texttt{predict\_mass} (rung $r_e=2$) and $c = \ell_0/\tau_0$:
this is an exact RS prediction.

\texttt{Lean: Masses.MassLaw.predict\_mass + Foundation.ConstantDerivations.c\_rs\_eq\_one}
\end{theorem}

\section{Derivation of Compton Scattering}

A photon of initial wavelength $\lambda$ scatters off a free electron
(at rest) through angle $\theta$.  Four-momentum conservation:
\begin{align}
  (E_\gamma + m_e c^2, \mathbf{p}_\gamma) &= (E_\gamma' + E_e, \mathbf{p}_\gamma' + \mathbf{p}_e).
\end{align}

Squaring the electron four-momentum and using $E_e^2 = p_e^2c^2 + m_e^2c^4$:

\begin{theorem}[Compton formula]\label{thm:compton}
\begin{equation}
  \Delta\lambda = \lambda' - \lambda = \frac{\hbar}{m_e c}(1 - \cos\theta)
  = \lambda_C(1-\cos\theta).
  \label{eq:compton}
\end{equation}
For $\theta = \pi/2$: $\Delta\lambda = \lambda_C = 2.426$ pm.
For $\theta = \pi$ (backscatter): $\Delta\lambda = 2\lambda_C = 4.852$ pm.
\end{theorem}

\begin{proof}
From four-momentum conservation:
$E_\gamma + m_e c^2 = E_\gamma' + E_e$, $\mathbf{p}_\gamma = \mathbf{p}_\gamma' + \mathbf{p}_e$.
Squaring: $(m_e c^2)^2 + 2m_e c^2(E_\gamma - E_\gamma') + 0 =
(m_e c^2)^2 + p_e^2 c^2$.
Also: $\mathbf{p}_e^2 = p_\gamma^2 + p_\gamma'^2 - 2p_\gamma p_\gamma'\cos\theta$.
Dividing by $E_\gamma E_\gamma'/(hc)^2$ gives eq.~\eqref{eq:compton}.
\end{proof}

\section{Derivation of the Photoelectric Effect}

\begin{theorem}[Photoelectric threshold]\label{thm:photo}
Light of frequency $\nu$ ejects electrons from a metal with maximum
kinetic energy:
\begin{equation}
  K_{\rm max} = h\nu - \phi_W = \hbar\omega - \phi_W,
\end{equation}
where $\phi_W$ is the work function.  Threshold: $h\nu_{\rm min} = \phi_W$.
No electrons are emitted for $\nu < \nu_{\rm min}$, regardless of intensity.

In RS: $E_{\rm photon} = \hbar\omega$ is the RS photon energy.  The work
function $\phi_W$ is the binding energy of the outermost electron, which
in RS is determined by the $\varphi$-ladder atomic energy structure.
\end{theorem}

\begin{remark}[Intensity independence]
In RS, the photon count (number of 8-tick ledger excitations) determines
the current, not intensity per se.  Each photon carries energy $\hbar\omega$.
A single photon with $\hbar\omega > \phi_W$ ejects exactly one electron.
This is enforced by the Pauli exclusion principle for the ejected electron
(Lean: \texttt{Foundation.SpinStatistics.pauli\_exclusion}).
\end{remark}

\section{Numerical Results}

\begin{center}
\begin{tabular}{llll}
\toprule
Quantity & RS Prediction & Experiment & Status \\
\midrule
$\lambda_C$ & $2.426 \times 10^{-12}$ m & $2.4263 \times 10^{-12}$ m & $< 0.01\%$ \\
Compton shift ($\theta=90°$) & $2.426$ pm & $2.426$ pm & \checkmark \\
Compton shift ($\theta=180°$) & $4.853$ pm & $4.852$ pm & \checkmark \\
Photoelectric: linear $K$ vs $\nu$ & $h\nu - \phi_W$ & Confirmed by Millikan & \checkmark \\
No emission below $\nu_{\rm min}$ & $\phi_W/h$ threshold & Confirmed & \checkmark \\
\bottomrule
\end{tabular}
\end{center}

\section{Lean Formalization Status}

\begin{center}
\begin{tabular}{llll}
\toprule
Claim & Lean theorem & Module & Status \\
\midrule
Electron mass $m_e$ & \texttt{predict\_mass} & \texttt{Masses.MassLaw} & Proved \\
Speed of light $c$ & \texttt{c\_rs\_eq\_one} & \texttt{Foundation.ConstantDerivations} & Proved \\
Photon spin-1 (bosonic) & \texttt{phase\_0\_is\_one} & \texttt{Foundation.EightTick} & Proved \\
Pauli exclusion (ejected e) & \texttt{pauli\_exclusion} & \texttt{Foundation.SpinStatistics} & Proved \\
Compton formula~\eqref{eq:compton} & Four-momentum conservation & --- & HYPOTHESIS$^\dagger$ \\
Work function from RS & --- & --- & HYPOTHESIS$^\ddagger$ \\
\bottomrule
\end{tabular}
\end{center}

$^\dagger$ \textbf{HYPOTHESIS}: Four-momentum conservation in Lean from RS
causal structure. Falsifier: Compton shift deviating from $\lambda_C(1-\cos\theta)$ at any angle.

$^\ddagger$ \textbf{HYPOTHESIS}: Work function $\phi_W$ derivable from RS $\varphi$-ladder atomic levels. Falsifier: systematic deviation of RS-predicted vs.\ measured $\phi_W$ for standard metals.

\section{Conclusion}

Compton scattering and the photoelectric effect both follow from the RS
photon ($E = \hbar\omega$, $p = \hbar\omega/c$) combined with the RS-derived
$m_e$ and $c$.  The Compton wavelength $\lambda_C = \hbar/(m_e c) = 2.426$ pm
is an exact zero-parameter RS prediction.

\begin{thebibliography}{99}
\bibitem{Einstein1905} A.~Einstein, ``On a heuristic viewpoint concerning the production and transformation of light,'' \textit{Ann.\ Phys.}\ \textbf{17}, 132 (1905).
\bibitem{Compton1923} A.~H.~Compton, ``A quantum theory of the scattering of X-rays by light elements,'' \textit{Phys.\ Rev.}\ \textbf{21}, 483 (1923).
\bibitem{Washburn2026a} J.~Washburn, ``The Algebra of Reality,'' \textit{Axioms} \textbf{15}(2), 90 (2026).
\end{thebibliography}

\end{document}
